% Generated draft: requires mathematical and source review.
\documentclass[11pt]{article}
\usepackage[margin=25mm]{geometry}
\usepackage{fontspec}
\setmainfont{Latin Modern Roman}
\usepackage{amsmath,amssymb,mathtools,unicode-math}
\setmathfont{Latin Modern Math}
\usepackage{enumitem,needspace,xurl,hyperref}
\setlength{\parindent}{0pt}
\setlength{\parskip}{5pt}
\setlength{\emergencystretch}{3em}
\newcommand{\N}{\mathbb{N}}
\newcommand{\Z}{\mathbb{Z}}
\newcommand{\Q}{\mathbb{Q}}
\newcommand{\R}{\mathbb{R}}
\newcommand{\C}{\mathbb{C}}
\newcommand{\F}{\mathbb{F}}
\newcommand{\A}{\mathbb{A}}
\newcommand{\PP}{\mathbb{P}}
\newcommand{\E}{\mathbb{E}}
\newcommand{\eps}{\varepsilon}
\newcommand{\dd}{\,\mathrm{d}}
\begin{document}
\section*{UnsolvedMath ID: 20002722}
\textbf{Source problem code:} catalogue entry
\subsection{Definitions and mathematical statement}
Problem 6.1 (H. Thomas, A. Zelevinsky) asks for a geometric polytopal realization of the cluster complex $\Delta_W$ for a noncrystallographic finite Coxeter group, and more strongly whether its known complete simplicial cluster fan is polytopal. Here $\Delta_W$ is the flag complex of compatible subsets of almost-positive roots $\Phi_{\geq-1}=\Phi^+\cup(-\Pi)$. Incidences require care: a simple polytope $P$ has cluster complex as the boundary complex of its simplicial polar $P^\circ$, not as its own face lattice. The normal-fan realization is the stronger geometric target.

The aid reports that the 2005 question is solved in the literature: Reading--Speyer identify the bipartite cluster fan linearly with a Cambrian fan; Hohlweg--Lange--Thomas realize finite-type Cambrian fans as normal fans of permutahedron-derived polytopes; and Felikson--Tumarkin--Yildirim give 2025 folding-section constructions including $H_3,H_4,I_2(m)$. These supplied status claims lack full references and must not be confused with a newly established theorem in this draft.

The aid additionally gives an explicit rank-two construction, valid for dihedral $I_2(m)$ with $m\geq3$. Let $u_1,\ldots,u_{m+2}$ be the almost-positive-root rays normalized to unit vectors in cyclic order. Let $\theta_j\in(0,\pi)$ be the counterclockwise angle from $u_j$ to $u_{j+1}$, using cyclic indices. Put
\[
Q_m=\operatorname{conv}\{u_j\},\qquad
P_m=Q_m^\circ=\{x:\langle u_j,x\rangle\leq1\ \forall j\}.
\]
The origin is interior, every $u_j$ is a vertex of $Q_m$, its boundary cycle is the cluster complex, and the normal fan of $P_m$ is exactly the given cluster fan. Its vertices and Euclidean metric data are
\[
v_j=\frac{u_j+u_{j+1}}{1+\cos\theta_j},\qquad
\operatorname{Area}(Q_m)=\frac12\sum_j\sin\theta_j,
\]
\[
\operatorname{Area}(P_m)=\sum_j\tan(\theta_j/2),\qquad
\operatorname{Perimeter}(P_m)=2\sum_j\tan(\theta_j/2),\qquad
\operatorname{Perimeter}(Q_m)=2\sum_j\sin(\theta_j/2).
\]
The normalization to unit rays fixes all metric constants; the combinatorial fan is unchanged by that normalization.

For positive support numbers $h_j$, set $P(h)=\{x:\langle u_j,x\rangle\leq h_j\ \forall j\}$ and let $v_j(h)$ be the unique intersection of the adjacent support lines $\langle u_j,x\rangle=h_j$ and $\langle u_{j+1},x\rangle=h_{j+1}$. The aid's necessary-and-sufficient finite stability test is
\[
\langle u_\ell,v_j(h)\rangle<h_\ell
\qquad\text{for every }j\text{ and every }\ell\notin\{j,j+1\}.
\]
These strict inequalities ensure that every adjacent intersection is a simple vertex with exactly the prescribed active normals, hence preserve the full normal fan. Non-strict feasibility alone may permit a coarser degenerate fan. This explicit rank-two polar model is separate from the cited higher-rank constructions and does not by itself build $H_3$ or $H_4$.
\subsection{Short English statement}
The source asks for noncrystallographic cluster polytopes, with the correct dual incidence and preferably the exact cluster normal fan. The aid reports general literature solutions and gives a direct dihedral model: take the convex hull of unit cluster rays and polarize it. Adjacent-line intersections provide explicit vertices and a strict finite test for preserving the fan under support changes.
\subsection{Sources}
\begin{itemize}
\item Local UnsolvedMath registry, permanent ID $20002722$, \texttt{catalogue entry}; \texttt{aim-probability-notes.json}, \texttt{notes[163]}, extracted source item 6.
\item American Institute of Mathematics, \emph{Braid groups, clusters and free probability}, Problem 6.1 (H. Thomas, A. Zelevinsky), \url{https://aimath.org/WWN/braidgroups/braidgroups.pdf}. Supplied verification repairs the split heading and subscript, and notes the preceding dual-associahedron/cyclohedron explanation.
\item Supplied literature attributions: Reading--Speyer; Hohlweg--Lange--Thomas; Felikson--Tumarkin--Yildirim (2025). Titles, theorem numbers, publication details and authoritative full citations are absent.
\item Research attempt 1: partial result, problem solved at run, rigor label ``proved,'' novelty confidence low. Registry classification solved in literature, status solved, statement status exact. The explicit rank-two metric and support test are supplied through the research aid.
\end{itemize}
\end{document}
